% Apresentacao Beamer -- Poco 861 Etapa 2 (campo Lapa)
% Versao com layout inspirado no template DRP (UENF--Equinor)
% Compilar (a partir desta pasta):
%   pdflatex poco861_etapa2_beamer_DRP.tex
%   pdflatex poco861_etapa2_beamer_DRP.tex
% Numeros alinhados a poco861_etapa2_dem_sc.tex (junho de 2026).

\documentclass[aspectratio=169,10pt]{beamer}

\usepackage[T1]{fontenc}
\usepackage[utf8]{inputenc}
\usepackage[brazil]{babel}
\usepackage{amsmath,amssymb}
\usepackage{booktabs}
\usepackage{array}
\usepackage{graphicx}
\usepackage{tikz}
\usetikzlibrary{positioning, arrows.meta, calc}
% DRP template uses Calibri; Carlito is the TeX metric-compatible substitute.
\usepackage[sfdefault]{carlito}
\usefonttheme{professionalfonts}

% Fundos DRP + figuras fig2_* do relatorio Etapa 2.
\graphicspath{%
  {./poco861_drp_assets/}%
  {./poco861_drp_assets/figures/}%
  {../figures/}%
  {../data/processed/dem_sc_runs/lab_validation/figures/}%
  {../data/processed/dem_sc_runs/multiscale_ab/figures/}%
  {../data/processed/dem_sc_runs/multiscale_nmr/figures/}%
  {../data/processed/dlis_861/figures/}%
}

% ------------------------------------------------------------------
% Paleta DRP
\definecolor{drpblue}{RGB}{29,55,113}
\definecolor{drporange}{RGB}{255,145,39}
\definecolor{drpyellow}{RGB}{250,238,0}
\definecolor{drpred}{RGB}{220,0,0}
\definecolor{drpgray}{RGB}{70,70,70}
\definecolor{po861blue}{RGB}{0,74,127}
\definecolor{po861green}{RGB}{0,128,96}

\usecolortheme{default}
\setbeamercolor{normal text}{fg=black,bg=white}
\setbeamercolor{structure}{fg=drpblue}
\setbeamercolor{frametitle}{fg=black,bg=}
\setbeamercolor{title}{fg=black,bg=}
\setbeamercolor{subtitle}{fg=drpgray,bg=}
\setbeamercolor{block title}{fg=white,bg=drpblue}
\setbeamercolor{block body}{fg=black,bg=drpblue!5}
\setbeamercolor{block title example}{fg=black,bg=drpyellow}
\setbeamercolor{block body example}{fg=black,bg=drpyellow!12}

\setbeamertemplate{navigation symbols}{}
\setbeamertemplate{blocks}[rounded][shadow=false]
\setbeamersize{text margin left=7mm, text margin right=7mm}

\setbeamerfont{normal text}{size=\small}
\setbeamerfont{itemize/enumerate body}{size=\small}
\setbeamerfont{itemize/enumerate subbody}{size=\footnotesize}
\setbeamerfont{block title}{size=\small,series=\bfseries}
\setbeamerfont{block body}{size=\footnotesize}
\setbeamerfont{frametitle}{size=\Large,series=\bfseries}
\setbeamerfont{title}{size=\Huge,series=\bfseries}
\setbeamerfont{subtitle}{size=\large}
\setbeamerfont{author}{size=\normalsize}
\setbeamerfont{institute}{size=\small}
\setbeamerfont{date}{size=\small}
\setbeamerfont{page number in head/foot}{size=\scriptsize,series=\mdseries}

\setlength{\parskip}{2pt}
\setlength{\itemsep}{2pt}
\setbeamertemplate{itemize item}{\raise0.2ex\hbox{\large\textbullet}}
\setbeamertemplate{itemize subitem}{\raise0.2ex\hbox{\small\textbullet}}
\setbeamercolor{itemize item}{fg=drpblue}
\setbeamercolor{itemize subitem}{fg=drpblue}

\usebackgroundtemplate{}

\setbeamertemplate{frametitle}{%
  \vspace*{0.18cm}%
  \hspace*{0.12cm}{\usebeamerfont{frametitle}\insertframetitle}\par%
  \vspace*{0.20cm}%
}

\newlength{\drpfooterheight}
\setlength{\drpfooterheight}{1.12cm}
\setbeamertemplate{footline}{%
  \begin{tikzpicture}[baseline=(base.base)]
    \node[inner sep=0pt,anchor=south west] (base) at (0,0)
      {
\includegraphics[width=\paperwidth,height=\drpfooterheight]{drp_footer.png}};
    \node[anchor=south east, font=\scriptsize, text=drpgray]
      at ([xshift=-0.27cm,yshift=0.13cm]base.south east) {\insertframenumber};
  \end{tikzpicture}%
}

\newcommand{\philab}{\ensuremath{\phi_{\mathrm{lab}}}}
\newcommand{\phiND}{\ensuremath{\phi_{\mathrm{ND}}}}
\newcommand{\phiS}{\ensuremath{\phi_{\mathrm{S}}}}
\newcommand{\Vp}{\ensuremath{V_{\mathrm{p}}}}
\newcommand{\Vs}{\ensuremath{V_{\mathrm{s}}}}
\newcommand{\VpVs}{\ensuremath{V_{\mathrm{p}}/V_{\mathrm{s}}}}

\newlength{\drpbodyheight}
\setlength{\drpbodyheight}{0.80\textheight}

\newcommand{\slidefig}[2][0.75]{%
  \includegraphics[width=\linewidth,height=#1\drpbodyheight,keepaspectratio]{#2}%
}
\newcommand{\slidecaption}[1]{%
  \par\vspace{0.05cm}%
  \begin{minipage}{\linewidth}%
    \centering\scriptsize #1%
  \end{minipage}%
}
\newcommand{\slidefigure}[3][0.72]{%
  \begin{minipage}[t]{\linewidth}%
    \centering
    \slidefig[#1]{#2}%
    \slidecaption{#3}%
  \end{minipage}%
}

\title{Poço~861 (campo Lapa): modelagem física DEM/SC}
\subtitle{Etapa~2 --- calibração por HFU e validação lab / perfil}
\author{Romulo Brito da Silva}
\institute{UENF / LENEP}
\date{Junho de 2026}

\begin{document}

% ------------------------------------------------------------------
{%
\usebackgroundtemplate{\includegraphics[width=\paperwidth,height=\paperheight]{drp_title_bg.png}}%
\begin{frame}[plain]
  \vspace*{0.40cm}
  \begin{center}
    {\huge\bfseries Poço~861\\[0.04cm]
    (campo Lapa)}\\[0.18cm]
    {\usebeamerfont{subtitle}\bfseries Modelagem física DEM/SC calibrada por HFU}\\[0.12cm]
    {\large Etapa~2 --- laboratório, perfil sónico e saturação}\\[0.22cm]
    {\usebeamerfont{author} Romulo Brito da Silva}\\[0.08cm]
    {\usebeamerfont{institute} UENF / LENEP}\\[0.08cm]
    {\usebeamerfont{date} Junho de 2026}\\[0.18cm]
    {\footnotesize
    Intervalo 5205,91--5233,72~m $\cdot$ 87 linhas de perfil\\
    10 plugs microCT $\cdot$ ROCKPHYS $\cdot$ DSI (DLIS) $\cdot$ CMR}\\[0.08cm]
    {\footnotesize Prova de conceito física \emph{local}}
  \end{center}
  \begin{tikzpicture}[remember picture,overlay]
    \node[anchor=south east, font=\scriptsize, text=drpgray]
      at ([xshift=-0.27cm,yshift=0.15cm]current page.south east) {\insertframenumber};
  \end{tikzpicture}
\end{frame}
}%
\usebackgroundtemplate{}

% ------------------------------------------------------------------
\begin{frame}{Roteiro}
  \footnotesize
  \begin{enumerate}
    \item Insumos da Etapa~1 que alimentam a física.
    \item Cadeia DEM/SC, papéis do microCT e do laboratório, premissas,
          partição de poros e ordem das inclusões.
    \item Duas escalas de erro e calibração por HFU.
    \item Resultados nos plugs (linha de base, LOO e multiescala CT).
    \item Extrapolação ao perfil: a porosidade de entrada versus o
          sónico de porosidade; depois o \Vp{} do DEM versus o DSI.
    \item Saturação (Gassmann), teste NMR e decisão da abordagem.
    \item Conclusões e caminho para a Etapa~3.
  \end{enumerate}
\end{frame}

% ------------------------------------------------------------------
\section{Contexto}

\begin{frame}{Da Etapa~1: insumos fixos para a física}
  \footnotesize
  A Etapa~1 definiu quais grandezas do perfil alimentam a física desta etapa.
  \vspace{0.08cm}
  \begin{columns}[T,onlytextwidth]
    \column{0.52\textwidth}
    \begin{block}{Entra na Etapa~2}
      \scriptsize
      \begin{itemize}
        \item \phiND{} como porosidade de entrada (default).
        \item HFU de laboratório (1--4) para segmentar.
        \item MicroCT: razão de aspecto e frações mineralógicas.
        \item Velocidades de laboratório (ROCKPHYS) e sónico DSI.
      \end{itemize}
    \end{block}
    \begin{exampleblock}{Fica de fora}
      \scriptsize
      \begin{itemize}
        \item ML direto perfil $\rightarrow$ \Vp{}/\Vs{}.
        \item FZI ou $k$ previstos só com wireline.
        \item MicroCT como feature extra de regressor.
      \end{itemize}
    \end{exampleblock}
    \column{0.46\textwidth}
    \centering
    \begin{tikzpicture}[
      node distance=0.85cm,
      flowbox/.style={
        draw, thick, rounded corners=4pt,
        align=center, text width=5.2cm,
        inner sep=5pt, font=\scriptsize
      },
      flowarr/.style={-{Stealth[length=2.8mm]}, thick, po861blue},
    ]
      \node[flowbox, fill=po861blue!18, draw=po861blue] (e1) {%
        \textbf{1. Diagnóstico ML}\\
        \phiND{}, HFU lab, microCT $\rightarrow$ física
      };
      \node[flowbox, fill=po861green!18, draw=po861green, below=of e1] (e2) {%
        \textbf{2. DEM/SC + Gassmann}\\
        calibração por HFU $\rightarrow$ \Vp{}, \Vs{}
      };
      \node[flowbox, fill=orange!18, draw=orange!70!black, below=of e2] (e3) {%
        \textbf{3. ML no resíduo}\\
        depois da física validada
      };
      \draw[flowarr] (e1.south) -- (e2.north);
      \draw[flowarr] (e2.south) -- (e3.north);
    \end{tikzpicture}
  \end{columns}
\end{frame}

% ------------------------------------------------------------------
\section{Dados e cadeia}

\begin{frame}{Fontes elásticas que fecham a Etapa~2}
  \footnotesize
  Além do perfil enriquecido da Etapa~1, entram velocidades de lab e de poço.
  \vspace{0.06cm}
  \centering
  \scriptsize
  \setlength{\tabcolsep}{2.5pt}
  \renewcommand{\arraystretch}{1.15}
  \begin{tabular}{@{}p{0.42\textwidth}p{0.34\textwidth}rc@{}}
    \toprule
    Fonte & Papel & $n$ & Escala \\
    \midrule
    Perfil integrado
      & porosidade neutrão-densidade (\phiND{}),
        unidade de fluxo hidrodinâmico (HFU),
        profundidade
      & 87 & $\sim$28~m \\
    Plugs de microtomografia
        computadorizada (microCT)
      & geometria de poros, porosidade de
        laboratório (\philab{})
      & 10 & pontual \\
    Planilha ROCKPHYS (Velocity)
      & velocidades compressional e
        cisalhante (\Vp{}, \Vs{}) de laboratório
      & 28 & plugs \\
    Sónico Dipole Shear Sonic Imager (DSI)
        em Digital Log Interchange
        Standard (DLIS)
      & lentidões compressional e
        cisalhante (DTCO/DTSM) \emph{in situ}
      & 182 & contínuo \\
    Combinable Magnetic Resonance (CMR) /
        Elemental Capture Spectroscopy (ECS)
        em DLIS
      & fluido livre (CMFF) e fluido ligado
        (BFV): frações macro/micro
      & 152 & contínuo \\
    \bottomrule
  \end{tabular}
  \vspace{0.10cm}
  \begin{block}{Distribuição de HFU no perfil}
    \scriptsize
    HFU1 $=42$ linhas, HFU2 $=29$, HFU3 $=10$, HFU4 $=6$.
    A HFU~4 não tem plug de laboratório; os parâmetros elásticos
    dessa unidade são a média das HFU~2 e~3 (aproximação, não medida).
  \end{block}
\end{frame}

\begin{frame}{Cadeia física: da matriz ao \Vp{} saturado}
  \footnotesize
  Ordem da cadeia: matriz $\rightarrow$ geometria $\rightarrow$ DEM seco
  $\rightarrow$ Gassmann. Os parâmetros elásticos são um conjunto
  \textbf{por HFU}.
  \vspace{0.06cm}
  \centering
  \begin{tikzpicture}[
    node distance=0.28cm and 0.18cm,
    box/.style={
      draw=po861blue, thick, rounded corners=2pt, fill=po861blue!10,
      minimum height=0.85cm, text width=2.15cm,
      align=center, font=\scriptsize, inner sep=2pt},
    arr/.style={-{Stealth[length=2.2mm]}, thick, po861blue},
  ]
    \node[box] (m) {Matriz VRH\\$K_m$, $G_m$};
    \node[box, right=of m] (a) {Aspecto $\alpha$\\(por HFU)};
    \node[box, right=of a] (d) {DEM/SC\\inclusões secas};
    \node[box, right=of d] (v) {\Vp{}, \Vs{}\\a seco};
    \node[box, right=of v, fill=po861green!12, draw=po861green] (g) {Gassmann\\$K_f\!\approx\!2{,}2$~GPa};
    \draw[arr] (m) -- (a);
    \draw[arr] (a) -- (d);
    \draw[arr] (d) -- (v);
    \draw[arr] (v) -- (g);
  \end{tikzpicture}
  \vspace{0.12cm}
  \begin{columns}[T,onlytextwidth]
    \column{0.48\textwidth}
    \begin{itemize}
      \item Entrada de porosidade: \phiND{} (perfil) ou \philab{} (plugs).
      \item $\alpha$ inicial: mediana microCT por HFU ($\bar{\alpha}$).
      \item Calibração inversa: $\alpha^\star$ e fator $s^\star$
            (comum a $K_m$ e $G_m$).
    \end{itemize}
    \column{0.48\textwidth}
    \begin{itemize}
      \item Inclusões secas ($K_i=G_i=0$) até $\phi$ alvo.
      \item Gassmann: saturação total ($S_w=1$).
      \item O sónico DSI entra só na validação.
    \end{itemize}
  \end{columns}
\end{frame}

\begin{frame}{Dois papéis no plug: imagem e laboratório}
  \footnotesize
  O microCT descreve \textbf{forma e minerais}. A porosidade de entrada do
  DEM no plug é a medida de laboratório \philab{}.
  \vspace{0.02cm}
  \centering
  \begin{tikzpicture}[
    src/.style={
      draw, thick, rounded corners=3pt, align=center,
      text width=3.5cm, inner sep=3pt, font=\scriptsize,
      minimum height=1.70cm
    },
    midbox/.style={
      draw=po861blue, thick, rounded corners=3pt, fill=po861blue!12,
      align=center, text width=5.4cm, inner sep=4pt, font=\scriptsize
    },
    resultbox/.style={
      draw=po861green, thick, rounded corners=3pt, fill=po861green!12,
      align=center, text width=5.4cm, inner sep=3pt, font=\scriptsize
    },
    arr/.style={-{Stealth[length=2.4mm]}, thick, po861blue},
    note/.style={font=\tiny, text=drpgray, align=center}
  ]
    \node[src, fill=orange!15, draw=orange!70!black] (ct) at (0,0) {%
      \textbf{MicroCT (imagem)}\\[0.04cm]
      $\alpha$ / AR; Solido1/2 (VRH)\\
      frações meso/micro\\
      {\tiny (só na multiescala)}
    };
    \node[src, fill=drpblue!8, draw=drpblue] (lab) at (4.8,0) {%
      \textbf{Laboratório}\\[0.04cm]
      $\phi = \philab{}$\\
      \Vp{}, \Vs{} de referência; HFU
    };
    \node[midbox] (dem) at (2.4,-1.95) {%
      \textbf{DEM / SC} --- inclusões secas
      $\rightarrow$ $K_{\mathrm{dry}}$, $G_{\mathrm{dry}}$
      $\rightarrow$ \Vp{} previsto
    };
    \node[resultbox] (cmp) at (2.4,-3.10) {%
      Compara com \Vp{} \textbf{lab}
    };
    \draw[arr] (ct.south) -- ++(0,-0.24) -| (dem.north);
    \draw[arr] (lab.south) -- ++(0,-0.24) -| (dem.north);
    \draw[arr] (dem) -- (cmp);
    \node[note] at (2.4,-3.75) {%
      No \textbf{perfil} (87 linhas): $\phi=\phiND{}$;
      do CT herdam-se parâmetros \textbf{por HFU}.
    };
  \end{tikzpicture}
  \vspace{0.02cm}
  \begin{exampleblock}{O que isso informa}
    \scriptsize
    A imagem informa \textbf{forma e composição};
    a porosidade do modelo vem do \textbf{laboratório}.
  \end{exampleblock}
\end{frame}

\begin{frame}{Premissas da modelagem física}
  \footnotesize
  Regras fixadas a priori. O sónico DSI só entra na validação.
  \vspace{0.06cm}
  \begin{columns}[T,onlytextwidth]
    \column{0.49\textwidth}
    \begin{block}{Petrofísica e textura}
      \scriptsize
      \begin{itemize}
        \item Porosidade de entrada: \phiND{} (perfil) / \philab{} (plugs).
        \item Segmentação por \textbf{HFU de laboratório}.
        \item Solido~1 $\rightarrow$ calcita; Solido~2 $\rightarrow$ dolomita
              (VRH; normalizado no CT).
        \item $\alpha$: achatamento do poro ($1=$ esfera; baixo $=$ mais
              achatado e mais mole). Ajuste só em $[0{,}15,\,0{,}95]$.
        \item HFU~4 sem plugs: usa a média das HFU~2 e~3
              (não há medida elástica nessa unidade).
      \end{itemize}
    \end{block}
    \column{0.49\textwidth}
    \begin{block}{Elástica e fluido}
      \scriptsize
      \begin{itemize}
        \item DEM Berryman a seco; SC só como controle.
        \item Inclusões vazias ($K_i = G_i = 0$).
        \item Gassmann \emph{depois} do arcabouço seco
              ($K_f \approx 2{,}2$~GPa, $S_w = 1$).
        \item $\alpha^\star$ e fator $s^\star$ (comum a $K_m$, $G_m$)
              ajustados nos \textbf{plugs} de laboratório.
        \item Limiar para adotar multiescala: ganho de MAPE $\geq 2$~p.p.
              no laboratório \emph{e} no perfil.
      \end{itemize}
    \end{block}
  \end{columns}
\end{frame}

\begin{frame}{Porosidade total e partição NMR}
  \footnotesize
  $\phiND{}$ diz \textbf{quanto} poro existe. O NMR diz \textbf{como}
  esse volume se reparte entre famílias (meso--macro e micro).
  \vspace{0.06cm}
  \begin{columns}[T,onlytextwidth]
    \column{0.48\textwidth}
    \begin{block}{1.\ $\phiND{}$ --- porosidade total}
      \scriptsize
      $\phi_{\mathrm{total}}(z)=\phiND{}(z)$ (perfil);
      nos plugs, $\philab{}$.
      Ex.: $\phiND{}=0{,}20$ $\Rightarrow$ o DEM introduz 20\% de poro.
    \end{block}
    \vspace{0.08cm}
    \begin{block}{2.\ NMR --- partição (CMFF/BFV)}
      \scriptsize
      $f_{\mathrm{meso}}=\dfrac{\mathrm{CMFF}}{\mathrm{CMFF}+\mathrm{BFV}}$,
      $f_{\mathrm{micro}}=\dfrac{\mathrm{BFV}}{\mathrm{CMFF}+\mathrm{BFV}}$,
      com $f_{\mathrm{meso}}+f_{\mathrm{micro}}=1$.
      No CT: frações GeoDict no lugar de CMFF/BFV.
    \end{block}
    \column{0.50\textwidth}
    \begin{exampleblock}{3.\ Combinação no DEM}
      \scriptsize
      $\phi_{\mathrm{meso}}=\phiND{}f_{\mathrm{meso}}$,
      $\phi_{\mathrm{micro}}=\phiND{}f_{\mathrm{micro}}$;
      a soma recupera $\phiND{}$.
      \vspace{0.06cm}
      Ex.: BFV$=0{,}06$, CMFF$=0{,}10$
      ($\phi_{\mathrm{NMR}}=0{,}16\neq 0{,}20$)
      $\Rightarrow$ $f_{\mathrm{micro}}=0{,}375$,
      $f_{\mathrm{meso}}=0{,}625$
      $\Rightarrow$ $\phi_{\mathrm{micro}}=0{,}075$,
      $\phi_{\mathrm{meso}}=0{,}125$.
      \vspace{0.06cm}
      Cada inclusão leva seu $\alpha$ (CT por HFU).
      Próximo slide: \emph{ordem} dessas inclusões.
    \end{exampleblock}
  \end{columns}
  \vspace{0.10cm}
  {\centering
  \tiny
  \textbf{Legenda.}
  CMFF $=$ FFI (fluido livre, curva CMR).
  BFV $=$ BVI (fluido ligado / irredutível, curva CMR).
  \phiND{} $=$ porosidade neutrão--densidade (perfil).
  \philab{} $=$ porosidade de laboratório (plugs).
  \par}
\end{frame}

\begin{frame}{Ordem das inclusões: volume decrescente (literatura)}
  \footnotesize
  Com $\phi_{\mathrm{inc}}=f_{\mathrm{pore}}\times\phi_{\mathrm{total}}$ já definido
  (slide anterior), o DEM sequencial é \emph{path-dependent}: a ordem
  precisa ser explícita.
  \vspace{0.06cm}
  \begin{columns}[T,onlytextwidth]
    \column{0.52\textwidth}
    \begin{block}{Convenção adotada}
      \scriptsize
      \begin{enumerate}
        \item Calcular $\phi_{\mathrm{inc}}$ para meso--macro e micro.
        \item Aplicar inclusões em ordem de $\phi_{\mathrm{inc}}$
              \textbf{decrescente} (maior volume $\rightarrow$ menor).
        \item Mesma regra da referência Equinor/Berryman do projeto.
      \end{enumerate}
    \end{block}
    \vspace{0.06cm}
    {\scriptsize
    A ordem segue o \textbf{volume} de cada família (Equinor/Berryman).
    Poros mais achatados (baixo $\alpha$) amolecem mais por unidade de
    volume; esse efeito entra em $\alpha$, não na ordem.}
    \vspace{0.08cm}
    \begin{block}{Boas práticas}
      \scriptsize
      Carbonatos: $\alpha$ meso/macro do microCT; micro via CT ou CMR.
      No poço~861 o modelo adotado permanece monoescala (ganho abaixo do limiar).
    \end{block}
    \column{0.46\textwidth}
    \centering
    \begin{tikzpicture}[
      node distance=0.42cm,
      box/.style={
        draw=po861blue, thick, rounded corners=3pt, fill=po861blue!10,
        text width=4.6cm, align=center, font=\scriptsize, inner sep=4pt
      },
      arr/.style={-{Stealth[length=2.4mm]}, thick, po861blue},
    ]
      \node[box] (m) {Matriz VRH ($K_m$, $G_m$)};
      \node[box, below=of m, fill=po861green!12, draw=po861green] (a) {%
        1.\ Inclusão de maior $\phi_{\mathrm{inc}}$\\
        (em geral meso--macro)
      };
      \node[box, below=of a, fill=orange!12, draw=orange!70!black] (b) {%
        2.\ Inclusão de menor $\phi_{\mathrm{inc}}$\\
        (em geral micro)
      };
      \node[box, below=of b] (o) {$K_{\mathrm{dry}}$, $G_{\mathrm{dry}}$
        $\rightarrow$ Gassmann};
      \draw[arr] (m) -- (a);
      \draw[arr] (a) -- (b);
      \draw[arr] (b) -- (o);
    \end{tikzpicture}
  \end{columns}
\end{frame}

% ------------------------------------------------------------------
\section{Protocolo}

\begin{frame}{Três escalas de validação --- três perguntas}
  \footnotesize
  Nenhuma escala sozinha aprova o modelo.
  \vspace{0.08cm}
  \centering
  \scriptsize
  \begin{tabular}{lll}
    \toprule
    Escala & Amostras & Pergunta \\
    \midrule
    Lab (linha de base) & 10 plugs CT & DEM cru vs \Vp{} lab? \\
    Lab in-sample / LOO & 10 plugs / 10 folds & Fecha o lab? Generaliza? \\
    Multiescala CT (A/B) & 10 plugs CT & DEM sequencial ganha vs M1? \\
    Perfil sónico DSI & 87 linhas & Extrapolação \emph{in situ}? \\
    Multiescala NMR & 87 linhas & Frações CMR melhoram M1? \\
    \bottomrule
  \end{tabular}
  \vspace{0.12cm}
  \begin{block}{O que as métricas informam}
    \scriptsize
    MAPE é o erro relativo típico em \Vp{}.
    RMSE penaliza os erros grandes e vem em km/s.
    Viés positivo: o modelo está mais rápido que o observado.
  \end{block}
\end{frame}

\begin{frame}{Como medimos a divergência: laboratório e poço}
  \footnotesize
  O DEM gera um \Vp{} \textbf{previsto}. A divergência é sempre
  previsto $-$ observado --- mas o observado muda de escala.
  \vspace{0.10cm}
  \centering
  \begin{tikzpicture}[
    box/.style={
      draw=po861blue, thick, rounded corners=3pt, fill=po861blue!12,
      align=center, text width=3.4cm, inner sep=5pt, font=\scriptsize
    },
    labbox/.style={
      draw=orange!70!black, thick, rounded corners=3pt, fill=orange!12,
      align=center, text width=4.0cm, inner sep=5pt, font=\scriptsize
    },
    wellbox/.style={
      draw=po861green, thick, rounded corners=3pt, fill=po861green!12,
      align=center, text width=4.0cm, inner sep=5pt, font=\scriptsize
    },
    arr/.style={-{Stealth[length=2.5mm]}, thick, po861blue},
    note/.style={font=\tiny, text=drpgray, align=center}
  ]
    \node[box] (dem) at (0,0) {%
      \textbf{\Vp{} previsto}\\
      DEM / SC\\
      (seco ou + Gassmann)
    };
    \node[labbox] (lab) at (5.2,1.15) {%
      \textbf{Observado: lab}\\
      \Vp{} ROCKPHYS (plugs)\\
      $\rightarrow$ MAPE, LOO
    };
    \node[wellbox] (well) at (5.2,-1.15) {%
      \textbf{Observado: poço}\\
      \Vp{} sónico DSI (perfil)\\
      $\rightarrow$ MAPE, viés
    };
    \draw[arr] (dem.east) -- ++(0.55,0) |- (lab.west)
      node[pos=0.75, above, font=\tiny, text=orange!70!black] {erro lab};
    \draw[arr] (dem.east) -- ++(0.55,0) |- (well.west)
      node[pos=0.75, below, font=\tiny, text=po861green] {erro poço};
    \node[note] at (2.6,-2.45) {%
      Cada escala tem o seu observado: plugs no laboratório, sónico no poço.
    };
  \end{tikzpicture}
  \vspace{0.08cm}
  \begin{exampleblock}{O que isso informa}
    \scriptsize
    São duas perguntas distintas: o modelo fecha o laboratório?
    Fecha o sónico no perfil?
  \end{exampleblock}
\end{frame}

\begin{frame}{Calibração por HFU: aprender no plug, aplicar no perfil}
  \footnotesize
  Parâmetros elásticos são ajustados no \textbf{lab} e só então levados às
  87 linhas --- sem reajuste ao sónico.
  \vspace{0.08cm}
  \centering
  \begin{tikzpicture}[
    box/.style={
      draw=po861blue, thick, rounded corners=3pt, fill=po861blue!12,
      align=center, text width=5.8cm, inner sep=5pt, font=\scriptsize
    },
    warn/.style={
      draw=drpred, thick, rounded corners=3pt, fill=red!6,
      align=center, text width=5.8cm, inner sep=4pt, font=\scriptsize
    },
    arr/.style={-{Stealth[length=2.5mm]}, thick, po861blue},
  ]
    \node[box] (a) {%
      \textbf{1. Aprender (10 plugs)}\\
      \philab{} + \Vp{} lab + microCT\\
      $\rightarrow$ ajusta $\alpha^\star$ e fator $s^\star$
      \textbf{por HFU}
    };
    \node[box, below=0.45cm of a, fill=po861green!12, draw=po861green] (b) {%
      \textbf{2. Aplicar (87 linhas)}\\
      $\phi=\phiND{}$ + HFU lab\\
      $\rightarrow$ usa $(\alpha^\star,s^\star)$ daquela HFU
    };
    \node[box, below=0.45cm of b, fill=orange!12, draw=orange!70!black] (c) {%
      \textbf{3. Validar no poço}\\
      compara \Vp{} DEM com sónico DSI\\
      (o DSI não entra no ajuste)
    };
    \node[warn, below=0.45cm of c] (d) {%
      HFU~4: sem plugs $\rightarrow$ usa a média das HFU~2 e~3\\
      (aproximação; confiança menor do que na HFU~2)
    };
    \draw[arr] (a) -- (b);
    \draw[arr] (b) -- (c);
    \draw[arr] (c) -- (d);
  \end{tikzpicture}
\end{frame}

\begin{frame}{Esquema por HFU: de $\bar{\alpha}$ a \Vp{} DEM}
  \footnotesize
  $\bar{\alpha}_i$ é só o ponto de partida (mediana CT). O modelo adotado usa
  $(\alpha_i^\star,s_i^\star)$ calibrados contra \Vp{} de laboratório.
  \vspace{0.04cm}
  \centering
  \begin{tikzpicture}[
    box/.style={
      draw=po861blue, thick, rounded corners=3pt, fill=po861blue!12,
      align=center, text width=10.6cm, inner sep=3.5pt, font=\scriptsize
    },
    arr/.style={-{Stealth[length=2.2mm]}, thick, po861blue},
  ]
    \node[box] (a) {%
      \textbf{1. Partida (microCT)}\\[0.03cm]
      $\displaystyle
        \bar{\alpha}_i
        =\mathrm{mediana}\{\alpha_{\mathrm{CT},j}:j\in\mathrm{HFU}_i\}
      $\quad;\quad
      matriz VRH $\rightarrow (K_m,G_m)_i$
    };
    \node[box, below=0.26cm of a, fill=po861green!12, draw=po861green] (b) {%
      \textbf{2. Calibração inversa (plugs, vs \Vp{} lab)}\\[0.03cm]
      $\displaystyle
        (\alpha_i^\star,s_i^\star)
        =\arg\min_{\alpha,s}\;
        \mathrm{erro}\!\big(
          V_p^{\mathrm{DEM}}(\phi_{\mathrm{lab}},\alpha,sK_m,sG_m),\;
          V_p^{\mathrm{lab}}
        \big)
      $\\[0.02cm]
      {\tiny partida: $\alpha=\bar{\alpha}_i$, $s=1$;\quad
       fator comum: $K_m^\star=s^\star K_m$, $G_m^\star=s^\star G_m$}
    };
    \node[box, below=0.26cm of b] (c) {%
      \textbf{3. Forward DEM (já calibrado)}\\[0.03cm]
      $\displaystyle
        (\alpha_i^\star,s_i^\star)
        \;\rightarrow\;
        \mathrm{DEM}(\phi,\alpha_i^\star,s_i^\star K_m,s_i^\star G_m)
        \;\rightarrow\;
        (K,G)_{\mathrm{dry}}
        \;\rightarrow\;
        (V_p,V_s)^{\mathrm{DEM}}
      $\\[0.02cm]
      {\tiny $\phi=\phi_{\mathrm{lab}}$ no plug;\quad
       $\phi=\phiND{}$ no perfil (87 linhas)}
    };
    \node[box, below=0.26cm of c, fill=orange!12, draw=orange!70!black] (d) {%
      \textbf{4. Validação (sem recalibrar)}\\[0.03cm]
      lab: $V_p^{\mathrm{DEM,seco}}$ vs $V_p^{\mathrm{lab}}$
      (ambos a seco)\\[0.02cm]
      poço: $V_p^{\mathrm{DEM}}+\mathrm{Gassmann}$ vs $V_p^{\mathrm{DSI}}(z)$
      (DSI é \emph{in situ} com fluido)
    };
    \draw[arr] (a) -- (b);
    \draw[arr] (b) -- (c);
    \draw[arr] (c) -- (d);
  \end{tikzpicture}
\end{frame}

% ------------------------------------------------------------------
\section{Calibração laboratorial}

\begin{frame}{Sem calibração: MAPE $\approx 26\%$, viés $+1{,}2$~km/s}
  \begin{columns}[T,onlytextwidth]
    \column{0.38\textwidth}
    \footnotesize
    DEM a seco com $\alpha$ do microCT e matriz VRH
    \textbf{superestima} \Vp{} de laboratório.
    \vspace{0.08cm}
    \begin{tabular}{lr}
      \toprule
      Métrica & Valor \\
      \midrule
      MAPE \Vp{} & 26,3\% \\
      RMSE & 1,37~km/s \\
      Viés & $+1{,}22$~km/s \\
      $n$ & 10 plugs \\
      \bottomrule
    \end{tabular}
    \par
    \vspace{0.10cm}
    {\scriptsize Sem um $\alpha$ por HFU o erro é grande e heterogêneo:
    cada unidade precisa da sua geometria efetiva, não de um valor único
    para o poço inteiro.}
    \column{0.60\textwidth}
    \slidefigure[0.86]{fig2_vp_dem_vs_lab}{DEM cru versus \Vp{} de laboratório
    (dez plugs CT, confinamento 22,1~MPa).}
  \end{columns}
\end{frame}

\begin{frame}{Calibração por HFU: MAPE in-sample cai para $\approx 9\%$}
  \begin{columns}[T,onlytextwidth]
    \column{0.42\textwidth}
    \footnotesize
    Ajuste de $\alpha^\star$ e do fator $s^\star$ (amolece $K_m$ e $G_m$
    em conjunto) \textbf{por HFU}, contra o \Vp{} de laboratório.
    \par
    \vspace{0.05cm}
    {\scriptsize
    $\alpha=1$ seria um poro esférico; valores baixos são poros
    achatados (mais moles). O ajuste só escolhe $\alpha$ em
    $[0{,}15,\,0{,}95]$.}
    \par
    \vspace{0.06cm}
    \scriptsize
    \setlength{\tabcolsep}{3pt}
    \begin{tabular}{cccc}
      \toprule
      HFU & $n$ & $\alpha_{\mathrm{cal}}$ & $\alpha_{\mathrm{CT}}$ \\
      \midrule
      1 & 4 & 0,95 & 0,55 \\
      2 & 4 & 0,18 & 0,55 \\
      3 & 2 & 0,95 & 0,61 \\
      4 & 0 & média 2--3 & --- \\
      \bottomrule
    \end{tabular}
    \par
    \vspace{0.08cm}
    {\scriptsize Nas HFU~1 e~3 o ótimo cola em 0,95: o modelo quer
    poros ainda mais redondos do que o intervalo permite.
    Na HFU~2 o ótimo é 0,18 (interior da faixa) --- unidade com
    melhor apoio experimental.}
    \column{0.56\textwidth}
    \slidefigure[0.72]{fig2_alpha_by_hfu}{$\alpha$ calibrado versus mediana
    microCT por HFU. Barras afastadas: a geometria da imagem, sozinha,
    não fecha o laboratório.}
  \end{columns}
\end{frame}

\begin{frame}{LOO: MAPE $\approx 15\%$ no laboratório}
  \begin{columns}[T,onlytextwidth]
    \column{0.40\textwidth}
    \footnotesize
    A cada rodada um plug fica de fora do ajuste.
    \vspace{0.08cm}
    \centering
    \scriptsize
    \begin{tabular}{lcc}
      \toprule
      Escopo & MAPE (\%) & RMSE \\
      \midrule
      In-sample & 9,1 & 0,52 \\
      LOO & 14,8 & 0,85 \\
      \bottomrule
    \end{tabular}
    \vspace{0.10cm}
    \begin{block}{O que a tabela informa}
      \scriptsize
      Os 9\% são o ajuste: os mesmos plugs calibram e são avaliados.
      No laboratório o número que vale é o LOO ($\approx 15\%$):
      a cada rodada um plug fica de fora, o modelo recalibra e só
      então se compara. Seis pontos a mais não surpreendem com
      $n=10$ --- em HFU de dois ou quatro plugs, tirar um é perder
      um quarto ou metade da amostra. O viés perto de zero indica
      que o modelo não está alto nem baixo de forma sistemática;
      o que resta é dispersão entre plugs.
    \end{block}
    \column{0.58\textwidth}
    \slidefigure[0.86]{fig2_vp_calib_before_after}{Antes versus depois da
    calibração (dez plugs).}
  \end{columns}
\end{frame}

\begin{frame}{Multiescala CT: manter monoescala}
  \footnotesize
  DEM sequencial meso-macro + micro (frações CT). Critério: ganho
  \emph{forward} $\geq 2$~p.p.\ com confirmação in-sample.
  \vspace{0.06cm}
  \begin{columns}[T,onlytextwidth]
    \column{0.40\textwidth}
    \scriptsize
    \begin{tabular}{llcc}
      \toprule
      Modelo & Escopo & MAPE & RMSE \\
      \midrule
      M1 & In-sample & 9,1 & 0,52 \\
      M2b forward & In-sample & 8,7 & 0,54 \\
      M1 & LOO & 14,8 & 0,85 \\
      M2b forward & LOO & 12,0 & 0,78 \\
      \bottomrule
    \end{tabular}
    \par
    \vspace{0.08cm}
    {\scriptsize Se cada plug tivesse o seu próprio
    $\alpha_{\mathrm{micro}}$ (teste diagnóstico, não operacional),
    o MAPE cairia até cerca de 3\%. No modo que pode ir ao perfil
    (mesmo $\alpha$ da HFU para todos os plugs da unidade), o ganho
    in-sample é só 0,4~p.p.\ e o RMSE piora.}
    \begin{exampleblock}{Decisão}
      \scriptsize Permanecer com um $\alpha$ efetivo por HFU nas
      87 linhas: a partição meso/micro do CT não paga o limiar de 2~p.p.
    \end{exampleblock}
    \column{0.58\textwidth}
    \slidefigure[0.84]{fig2_multiscale_ab_crossplot}{M1 versus M2b
    \emph{forward} (in-sample).}
  \end{columns}
\end{frame}

% ------------------------------------------------------------------
\section{Perfil e sónico}

\begin{frame}{A porosidade de entrada alinha-se ao sónico; \Vp{} ainda não}
  \begin{columns}[T,onlytextwidth]
    \column{0.42\textwidth}
    \footnotesize
    Os parâmetros calibrados nos plugs são aplicados às 87 linhas
    (\phiND{} + HFU de laboratório).
    \par
    \vspace{0.08cm}
    A figura ao lado compara \textbf{duas porosidades}, não duas
    velocidades: \phiND{} (entrada do DEM) versus \phiS{}
    (porosidade inferida do sónico).
    \par
    \vspace{0.08cm}
    {\scriptsize
    Pontos junto à diagonal e $r \approx 0{,}94$ mostram que a
    \textbf{entrada petrofísica} está consistente com o perfil.
    Isso \textbf{não} prova que o \Vp{} do DEM está certo ---
    essa pergunta é o confronto com o DSI, no próximo slide.}
    \par
    \vspace{0.10cm}
    \scriptsize
    \setlength{\tabcolsep}{3pt}
    \begin{tabular}{lr}
      \toprule
      Grandeza (DEM, 87 linhas) & Valor \\
      \midrule
      \Vp{} médio a seco & 5,20~km/s \\
      \VpVs{} médio & 1,82 \\
      $r$(\phiND{}, \phiS{}) & 0,94 \\
      \bottomrule
    \end{tabular}
    \column{0.56\textwidth}
    \slidefigure[0.80]{fig2_phi_nd_vs_phi_sonic}{\phiND{} versus porosidade
    sónica \phiS{} ($n=87$). Acordo da \textbf{entrada}, não do \Vp{}.}
  \end{columns}
\end{frame}

\begin{frame}{Contra o DSI a seco: MAPE $\approx 9\%$, viés $+0{,}30$~km/s}
  \footnotesize
  87/87 profundidades casadas.
  {\scriptsize\quad
  MAPE 8,7\% $\cdot$
  RMSE 0,54~km/s $\cdot$
  viés $+0{,}30$~km/s $\cdot$
  Pearson $r=0{,}34$ $\cdot$
  \Vp{} DEM 5,20 vs sónico 4,90~km/s.}
  \vspace{0.06cm}
  \begin{columns}[T,onlytextwidth]
    \column{0.48\textwidth}
    \slidefigure[0.72]{fig2_vp_dem_vs_sonic_depth}{Trilha: DEM seco versus
    sónico DSI.}
    \column{0.50\textwidth}
    \slidefigure[0.72]{fig2_vp_dem_vs_sonic_scatter}{Dispersograma 1:1 por HFU
    (monoescala seco).}
  \end{columns}
  \vspace{0.04cm}
  {\scriptsize O DEM está, em média, 0,30~km/s mais rápido que o DSI.
  Esse deslocamento sistemático responde por cerca de três quartos do
  erro absoluto: não é só dispersão pontual --- é um viés que a Etapa~3
  vai corrigir.}
\end{frame}

\begin{frame}{Por HFU: os maiores erros concentram-se onde há menos suporte}
  \footnotesize
  O erro global esconde heterogeneidade forte entre unidades.
  \vspace{0.08cm}
  \begin{columns}[T,onlytextwidth]
    \column{0.42\textwidth}
    {\centering
    \scriptsize
    \setlength{\tabcolsep}{2.5pt}
    \begin{tabular}{ccccc}
      \toprule
      HFU & $n$ & Plugs & $\alpha_{\mathrm{cal}}$ & MAPE (\%) \\
      \midrule
      1 & 42 & 4 & 0,95$^\star$ & 8,0 \\
      2 & 29 & 4 & 0,18 & \textbf{3,3} \\
      3 & 10 & 2 & 0,95$^\star$ & 13,6 \\
      4 & 6 & 0 & média 2--3 & 30,4 \\
      \bottomrule
    \end{tabular}
    \par}
    \vspace{0.04cm}
    {\scriptsize $^\star$ $\alpha$ no máximo permitido (0,95): o ajuste
    pediu poros ainda mais esféricos do que $[0{,}15,\,0{,}95]$ admite.}
    \vspace{0.05cm}
    \begin{block}{O que a tabela informa}
      \scriptsize
      HFU~1 e HFU~2 têm \textbf{as mesmas} 4 plugs, mas MAPE 8,0\%
      contra 3,3\%. O que muda é a geometria efetiva: na HFU~2
      $\alpha=0{,}18$ (poro achatado, interior da faixa) e o DSI fecha
      bem; na HFU~1 $\alpha$ cola em 0,95 (quase esfera) e o erro sobe.
      \par
      A HFU~4 não tem plug: herda a média das HFU~2 e~3.
      MAPE de 30\% --- é a unidade em que novos plugs mais reduziriam
      a incerteza.
    \end{block}
    \column{0.56\textwidth}
    \slidefigure[0.72]{fig2_vp_error_vs_depth}{Erro de \Vp{} versus
    profundidade (DEM seco vs DSI). Os picos coincidem com a HFU~4.}
  \end{columns}
\end{frame}

% ------------------------------------------------------------------
\section{Gassmann e multiescala NMR}

\begin{frame}{Gassmann leva o MAPE a $\approx 7\%$, viés persiste}
  \begin{columns}[T,onlytextwidth]
    \column{0.36\textwidth}
    \footnotesize
    Fluido default: $K_f \approx 2{,}2$~GPa, $S_w=1$.
    \vspace{0.06cm}
    \centering
    \scriptsize
    \begin{tabular}{lccc}
      \toprule
      Condição & MAPE & RMSE & Viés \\
      \midrule
      DEM seco & 8,7 & 0,54 & $+0{,}30$ \\
      + Gassmann & 7,3 & 0,48 & $+0{,}27$ \\
      \bottomrule
    \end{tabular}
    \par
    \vspace{0.10cm}
    {\scriptsize Com porosidade média de cerca de 11\% e matriz rígida,
    encher o poro aumenta $K$ (aceleraria \Vp{}) e a densidade
    (desacelera). Os dois efeitos quase se cancelam: o MAPE cai pouco
    (8,7\% $\rightarrow$ 7,3\%) e o viés permanece. O que resta é o
    residual da Etapa~3.}
    \column{0.62\textwidth}
    \slidefigure[0.86]{fig2_vp_dry_gassmann_sonic_depth}{Sónico, DEM seco e
    DEM + Gassmann ao longo da profundidade.}
  \end{columns}
\end{frame}

\begin{frame}{NMR ponto a ponto: o \Vp{} médio acerta, a dispersão aumenta}
  \footnotesize
  As frações CMR (CMFF/BFV) variam em cada profundidade e podem
  substituir a mediana CT por HFU.
  \vspace{0.04cm}
  \begin{columns}[T,onlytextwidth]
    \column{0.40\textwidth}
    \scriptsize
    \begin{itemize}
      \item Merge CMR + perfil: \textbf{87/87}.
      \item M3 = DEM sequencial + Gassmann.
      \item $\alpha_{\mathrm{meso/micro}}$ fixos por HFU (teste CT).
    \end{itemize}
    \vspace{0.04cm}
    \centering
    \begin{tabular}{lccc}
      \toprule
      Modelo & MAPE & Viés & $r$ \\
      \midrule
      M1 Gassmann & 7,3 & $+0{,}27$ & $+0{,}42$ \\
      M3 NMR & 8,5 & $+0{,}04$ & $+0{,}45$ \\
      \bottomrule
    \end{tabular}
    \par
    \vspace{0.06cm}
    {\scriptsize M3 quase zera o viés
    ($+0{,}27 \rightarrow +0{,}04$~km/s): o \Vp{} médio cai no
    lugar. O MAPE, porém, sobe de 7,3\% para 8,5\% --- o acerto
    médio melhora e o ponto a ponto piora.
    CMR vs CT nos plugs: $r \approx -0{,}30$.}
    \begin{exampleblock}{Decisão}
      \scriptsize Manter monoescala Gassmann nas 87 linhas.
    \end{exampleblock}
    \column{0.58\textwidth}
    \slidefigure[0.82]{fig2_multiscale_nmr_crossplot_vp}{M1 versus M3 no
    plano 1:1 (sónico no eixo $y$).}
  \end{columns}
\end{frame}

\begin{frame}{O que o teste NMR isolou}
  \footnotesize
  Comparação controlada: só as frações de poro passam a variar com a
  profundidade (CMFF/BFV). $\phiND{}$, $\alpha$ e a matriz permanecem
  como no M1.
  \vspace{0.08cm}
  \centering
  \begin{tikzpicture}[
    knob/.style={
      draw, thick, rounded corners=2pt, align=center,
      text width=2.15cm, inner sep=3pt, font=\scriptsize, minimum height=0.95cm
    },
    arr/.style={-{Stealth[length=2.2mm]}, thick, po861blue},
  ]
    \node[font=\small\bfseries] at (0,1.55) {M1 monoescala};
    \node[knob, fill=po861blue!10, draw=po861blue] (m1a) at (-3.3,0.55) {$\phi=\phiND{}$};
    \node[knob, fill=po861blue!10, draw=po861blue] (m1b) at (-1.1,0.55) {$\alpha$ efetivo\\por HFU};
    \node[knob, fill=po861blue!10, draw=po861blue] (m1c) at (1.1,0.55) {matriz\\$K,G$};
    \node[knob, fill=po861blue!10, draw=po861blue] (m1d) at (3.3,0.55) {Gassmann};

    \node[font=\small\bfseries] at (0,-0.55) {M3 NMR multiescala};
    \node[knob, fill=po861blue!10, draw=po861blue] (m3a) at (-3.3,-1.55) {$\phi=\phiND{}$\\{\tiny (igual)}};
    \node[knob, fill=drpyellow!40, draw=orange!80!black] (m3b) at (-1.1,-1.55) {%
      \textbf{frações NMR}\\{\tiny CMFF/BFV $\star$}
    };
    \node[knob, fill=gray!15, draw=gray] (m3c) at (1.1,-1.55) {%
      AR meso/micro\\{\tiny \textbf{fixo} por HFU}
    };
    \node[knob, fill=po861blue!10, draw=po861blue] (m3d) at (3.3,-1.55) {matriz +\\Gassmann};

    \draw[arr, orange!80!black] (m1b.south) -- (m3b.north)
      node[midway, right, font=\tiny, text=orange!80!black] {mudou};
    \draw[arr, gray] (m1b.south) -- ++(0,-0.35) -| (m3c.north);
  \end{tikzpicture}
  \vspace{0.04cm}
  \begin{columns}[T,onlytextwidth]
    \column{0.48\textwidth}
    \begin{block}{O que o NMR trouxe}
      \scriptsize
      Quanto é macro vs micro \textbf{em cada profundidade}.
    \end{block}
    \column{0.48\textwidth}
    \begin{exampleblock}{O que permaneceu constante}
      \scriptsize
      Quão achatado é o poro ($\alpha$) --- ainda um valor por HFU.
    \end{exampleblock}
  \end{columns}
  \vspace{0.04cm}
  {\scriptsize $\alpha$ controla a moleza (poro achatado amolece mais).
  Trocar só as frações desloca o \Vp{} médio (o viés cai), mas o MAPE
  sobe 1,2~p.p.\ porque a forma do poro continua constante por HFU.}
\end{frame}

\begin{frame}{Decisão multiescala: limiar de 2~p.p.}
  \footnotesize
  Critério definido a priori: adotar multiescala se o ganho de MAPE for
  de pelo menos 2~p.p.\ e aparecer no laboratório e no perfil.
  \vspace{0.08cm}
  \centering
  \begin{tikzpicture}[
    box/.style={
      draw=po861blue, thick, rounded corners=3pt, fill=po861blue!12,
      align=center, text width=2.55cm, inner sep=3pt, font=\scriptsize
    },
    q/.style={
      draw=orange!70!black, thick, rounded corners=3pt, fill=orange!12,
      align=center, text width=3.2cm, inner sep=4pt, font=\scriptsize
    },
    yes/.style={
      draw=po861green, thick, rounded corners=3pt, fill=po861green!12,
      align=center, text width=2.6cm, inner sep=3pt, font=\scriptsize
    },
    no/.style={
      draw=drpred, thick, rounded corners=3pt, fill=red!6,
      align=center, text width=3.0cm, inner sep=3pt, font=\scriptsize
    },
    arr/.style={-{Stealth[length=2.4mm]}, thick, po861blue},
  ]
    \node[box] (m1) at (0,1.35) {M1 monoescala\\(baseline)};
    \node[box] (m2) at (3.3,1.35) {M2b\\multiescala CT};
    \node[box] (m3) at (6.6,1.35) {M3\\multiescala NMR};
    \node[q] (q) at (3.3,0.0) {Ganho MAPE $\geq 2$~p.p.\ vs M1?};
    \node[yes] (y) at (0.6,-1.45) {Sim $\rightarrow$\\adotar multi};
    \node[no] (n) at (5.6,-1.45) {Não $\rightarrow$\\
      \textbf{manter mono}\\{\tiny (caso 861)}};
    \draw[arr] (m1) -- (q);
    \draw[arr] (m2) -- (q);
    \draw[arr] (m3) -- (q);
    \draw[arr] (q) -- (y);
    \draw[arr] (q) -- (n);
  \end{tikzpicture}
  \vspace{0.08cm}
  \begin{exampleblock}{No 861}
    \scriptsize
    Multiescala CT: forward ganha 0,4~p.p.\ in-sample.
    Multiescala NMR: MAPE piora 1,2~p.p.
    Abordagem adotada: monoescala Gassmann + $\alpha$ efetivo por HFU.
  \end{exampleblock}
\end{frame}

\begin{frame}{Síntese das validações de \Vp{}}
  \footnotesize
  Dois regimes de erro: dispersão no lab, viés no perfil.
  \vspace{0.08cm}
  \centering
  \scriptsize
  \begin{tabular}{lcccc}
    \toprule
    Escala & MAPE (\%) & RMSE & Viés & $n$ \\
    \midrule
    Lab, sem calibração & 26,3 & 1,37 & $+1{,}22$ & 10 \\
    Lab, in-sample & 9,1 & 0,52 & $-0{,}01$ & 10 \\
    Lab, LOO & 14,8 & 0,85 & $-0{,}07$ & 10 \\
    Perfil DSI (seco) & 8,7 & 0,54 & $+0{,}30$ & 87 \\
    Perfil DSI (Gassmann) & 7,3 & 0,48 & $+0{,}27$ & 87 \\
    Perfil DSI (NMR M3) & 8,5 & 0,47 & $+0{,}04$ & 87 \\
    \bottomrule
  \end{tabular}
  \vspace{0.12cm}
  \begin{block}{O que a tabela informa}
    \scriptsize
    No laboratório o erro é dispersão (viés nulo em dez plugs).
    No perfil o erro é sobretudo um deslocamento positivo --- alvo da
    correção residual na Etapa~3.
  \end{block}
\end{frame}

% ------------------------------------------------------------------
\section{Conclusões}

\begin{frame}{Conclusões e recomendações}
  \footnotesize
  \begin{enumerate}
    \item \textbf{Viabilidade local:} DEM/SC por HFU fecha o perfil
          ($\approx 9\%$ seco, $\approx 7\%$ com Gassmann) e o lab sob a
          métrica mais severa (LOO $\approx 15\%$).
    \item \textbf{Uso por HFU:} o acordo com o DSI não é uniforme.
          HFU~2 ($\alpha=0{,}18$, 4 plugs): 3\%.
          HFU~1 ($\alpha$ no máximo 0,95, 4 plugs): 8\%.
          HFU~3 (2 plugs, $\alpha$ no máximo): 14\%.
          HFU~4 (sem plug, parâmetros herdados): 30\% ---
          unidade em que novos dados mais ajudariam.
    \item \textbf{Entrada:} \phiND{} no perfil; microCT na geometria e nos
          minerais da física.
    \item \textbf{Multiescala CT e NMR:} o ganho fica abaixo de 2~p.p.\ ---
          modelo monoescala com Gassmann.
    \item \textbf{Gassmann:} MAPE 8,7\% $\rightarrow$ 7,3\%, viés
          $+0{,}30 \rightarrow +0{,}27$~km/s; reduz o afastamento e deixa
          um residual estruturado.
    \item \textbf{Etapa~3:} ML sobre o resíduo
          (sónico $-$ física), com prioridade nas HFU~3 e~4.
  \end{enumerate}
\end{frame}

\begin{frame}{Caminho: física validada $\rightarrow$ ML residual}
  \vspace{-0.08cm}
  \begin{columns}[T,onlytextwidth]
    \column{0.50\textwidth}
    \footnotesize
    \textbf{Resultado da Etapa~2}
    \begin{itemize}
      \item \Vp{}/\Vs{} teóricos por HFU (seco + Gassmann).
      \item Parâmetros calibrados rastreáveis.
      \item Decisão documentada: monoescala.
    \end{itemize}
    \textbf{Próximo (Etapa~3)}
    \begin{itemize}
      \item Alvo: resíduo $= \Vp^{\mathrm{sonico}} - \Vp^{\mathrm{fisica}}$.
      \item Validação: blocos de profundidade (protocolo Etapa~1).
      \item Híbrido: física + correção residual pontual.
    \end{itemize}
    \column{0.48\textwidth}
    \centering
    \begin{tikzpicture}[
      node distance=0.90cm,
      flowbox/.style={
        draw, thick, rounded corners=4pt,
        align=center, text width=5.4cm,
        inner sep=5pt, font=\scriptsize
      },
      flowarr/.style={-{Stealth[length=2.8mm]}, thick, po861blue},
      arrlabel/.style={
        midway, fill=white, inner sep=2pt,
        font=\tiny, text=po861blue, align=center
      }
    ]
      \node[flowbox, fill=po861blue!18, draw=po861blue] (e1) {%
        \textbf{1. Diagnóstico ML}\\
        Limites do perfil
      };
      \node[flowbox, fill=po861green!18, draw=po861green, below=of e1] (e2) {%
        \textbf{2. DEM/SC + Gassmann}\\
        \textbf{concluída nesta apresentação}
      };
      \node[flowbox, fill=orange!18, draw=orange!70!black, below=of e2] (e3) {%
        \textbf{3. Física + ML residual}\\
        Ajuste fino de \Vp{} (OOF)
      };
      \draw[flowarr] (e1.south) --
        node[arrlabel] {$\phi$, HFU} (e2.north);
      \draw[flowarr] (e2.south) --
        node[arrlabel] {resíduo} (e3.north);
    \end{tikzpicture}
  \end{columns}
  \vspace{0.12cm}
  \begin{center}
    {\small Obrigado.}
  \end{center}
\end{frame}

\end{document}
